\section*{Abstract}\addcontentsline{toc}{section}{Abstract}

Holistic Mind is a mobile wellness application that combines daily check-ins, personalised practice recommendations, reflective journaling, and structured learning content. Originally proposed as trauma-informed support between professional appointments for adults experiencing anxiety, ADHD, or trauma-related difficulties, the application retains a voluntary, non-diagnostic wellness scope. The project investigates how a small, explainable recommendation system can support a daily self-care routine while respecting explicitly declared comfort preferences and keeping journal content private. The artefact consists of an Expo React Native application, a Node.js and Express API, PostgreSQL, an administration dashboard, object storage, and a Python recommendation service. Its hybrid engine combines suitability rules, MiniLM semantic similarity, and, when sufficient interactions exist, collaborative evidence. Journal entries are encrypted on the device and are excluded from server-side recommendation inputs. A saved journal-free development benchmark covers 30 authored scenarios and a fixed 14-exercise catalogue. The hybrid approach achieved Precision@4 of 77.50\% and NDCG@4 of 0.7387, compared with 73.33\% and 0.7244 for rules with eligibility filtering. However, hybrid recommendations also included 3.33\% negative-labelled selections, showing that ranking improvements do not establish universal suitability. A separate eight-scenario comfort suite recorded no explicit exclusion violations. Documented simulator and backend checks support functional feasibility, while independent label review, physical-device testing, user acceptance research, and clinical effectiveness remain unverified. The contribution is an integrated, privacy-conscious software artefact with transparent evaluation and identifiable limitations.
